% =====================================================================
%  Macros de figures
% =====================================================================
%  Toutes les figures du rapport passent par ces macros : la mise en
%  forme (centrage, largeur, legende, espacement) reste ainsi identique
%  d'un bout a l'autre du document.
%
%  Le guide impose que chaque figure porte une legende ET soit citee
%  dans le texte. Chaque macro prend donc un label, et le texte doit
%  contenir un \ref correspondant.
% =====================================================================

\graphicspath{{figures/}{figures/diagrammes/}{figures/logos/}%
              {figures/captures/}{figures/planning/}}

% ---------------------------------------------------------------------
%  \figReelle{fichier}{legende}{label}{largeur}
%  Figure ordinaire. La largeur s'exprime en fraction de \textwidth.
% ---------------------------------------------------------------------
%  Placement « htbp » et non « H » : un diagramme agrandi ne tient
%  souvent pas dans le bas de page ou il tombe. Ancre de force, il y
%  laisse un vide et se reporte a la page suivante ; laisse flottant, il
%  se pose en tete de page et le texte comble l'intervalle. Le gain
%  mesure sur ce document est de sept pages, sans perte de proximite
%  puisque chaque diagramme est cite dans le texte par son numero.
\newcommand{\figReelle}[4]{%
  \begin{figure}[htbp]
    \centering
    \includegraphics[width=#4\textwidth,keepaspectratio]{#1}
    \caption{#2}
    \label{#3}
  \end{figure}%
}

% ---------------------------------------------------------------------
%  \capture{fichier}{legende}{label}
%  Capture d'ecran : large, encadree d'un filet fin pour la detacher du
%  fond de page. Le regroupement de quatre ou cinq captures sous une
%  seule figure a ete abandonne : reduites a moins de la moitie de la
%  largeur du texte, elles devenaient illisibles a l'impression.
%
%  LARGEUR UNIQUE. Les trente-sept captures du rapport ont ete prises
%  dans la meme fenetre : elles partagent le meme rapport, 2,07. Une
%  largeur unique suffit donc, et c'est elle qui donne a une page
%  portant deux captures son aplomb. Toute exception se verrait.
% ---------------------------------------------------------------------
\newcommand{\largeurCapture}{0.90}

\newcommand{\capture}[3]{%
  \begin{figure}[H]
    \centering
    \fbox{\includegraphics[width=\largeurCapture\textwidth,keepaspectratio]{#1}}
    \caption{#2}
    \label{#3}
  \end{figure}%
}

%  Conservee pour compatibilite : elle composait autrefois plus etroit.
%  Elle ne le fait plus. La seule capture qui l'employait, celle de
%  l'application chauffeur, mesure 1887 x 912 pixels comme les autres :
%  la reduire a 0,60 la rendait simplement plus petite que ses voisines,
%  sans raison.
\newcommand{\captureEtroite}[3]{\capture{#1}{#2}{#3}}

% ---------------------------------------------------------------------
%  \figDeux / \figTrois : plusieurs images sous UNE seule legende.
%  Reservees aux DIAGRAMMES qui se lisent ensemble. Ne plus les employer
%  pour des captures d'ecran.
% ---------------------------------------------------------------------
\newcommand{\figDeux}[5]{%
  \begin{figure}[H]
    \centering
    \includegraphics[width=#5\textwidth,keepaspectratio]{#1}\\[4pt]
    \includegraphics[width=#5\textwidth,keepaspectratio]{#2}
    \caption{#3}
    \label{#4}
  \end{figure}%
}

\newcommand{\figTrois}[6]{%
  \begin{figure}[H]
    \centering
    \includegraphics[width=#6\textwidth,keepaspectratio]{#1}\\[3pt]
    \includegraphics[width=#6\textwidth,keepaspectratio]{#2}\\[3pt]
    \includegraphics[width=#6\textwidth,keepaspectratio]{#3}
    \caption{#4}
    \label{#5}
  \end{figure}%
}

% ---------------------------------------------------------------------
%  \figPaysage{fichier}{legende}{label}
%  Diagramme trop large pour la justification : page en paysage.
% ---------------------------------------------------------------------
\newcommand{\figPaysage}[3]{%
  \begin{landscape}
    \begin{figure}[H]
      \centering
      \includegraphics[width=0.95\linewidth,height=0.82\textheight,%
                       keepaspectratio]{#1}
      \caption{#2}
      \label{#3}
    \end{figure}
  \end{landscape}%
}

% =====================================================================
%  Emplacements reserves
% =====================================================================
%  Une figure encore absente occupe deja sa place, avec sa VRAIE
%  legende et son VRAI label : elle est donc numerotee, citee dans le
%  texte et listee dans la table des figures. Le jour ou l'image
%  existe, seule la commande change : \figAttendue devient \figReelle.
% =====================================================================
\definecolor{cadreAttente}{HTML}{B0562A}
\definecolor{fondAttente}{HTML}{FDF3EC}

\newcommand{\figAttendue}[3]{%
  \begin{figure}[H]
    \centering
    \begin{tcolorbox}[width=0.86\textwidth,height=5.4cm,
      colback=fondAttente,colframe=cadreAttente,boxrule=1pt,
      borderline={0.5pt}{3pt}{cadreAttente,dashed},
      halign=center,valign=center,sharp corners]
      {\sffamily\bfseries\color{cadreAttente}FIGURE A FOURNIR}\\[6pt]
      {\sffamily\small Fichier attendu :}\\[2pt]
      {\ttfamily\small #1}
    \end{tcolorbox}
    \caption{#2}
    \label{#3}
  \end{figure}%
}

%  Variante pour les captures d'ecran : indique en plus l'ecran a
%  photographier, pour que la reprise soit immediate.
\newcommand{\captureAttendue}[4]{%
  \begin{figure}[H]
    \centering
    \begin{tcolorbox}[width=0.86\textwidth,height=5.4cm,
      colback=fondAttente,colframe=cadreAttente,boxrule=1pt,
      borderline={0.5pt}{3pt}{cadreAttente,dashed},
      halign=center,valign=center,sharp corners]
      {\sffamily\bfseries\color{cadreAttente}CAPTURE A FOURNIR}\\[6pt]
      {\sffamily\small Ecran :} {\sffamily\itshape #4}\\[4pt]
      {\sffamily\small Fichier attendu :}\\[2pt]
      {\ttfamily\small #1}
    \end{tcolorbox}
    \caption{#2}
    \label{#3}
  \end{figure}%
}

% =====================================================================
%  Encadres de contenu
% =====================================================================
\newtcolorbox{definitionbox}[1]{%
  breakable, colback=grisFond, colframe=grisFilet, boxrule=0.6pt,
  left=8pt, right=8pt, top=6pt, bottom=6pt, sharp corners,
  title={\sffamily\bfseries #1}, coltitle=bleuTitre,
  colbacktitle=grisFond, titlerule=0pt}

% Mise en valeur imposee par le guide : les exemples en italique, les
% definitions encadrees, les elements essentiels en gras.
\newcommand{\exemple}[1]{\textit{#1}}
\newcommand{\essentiel}[1]{\textbf{#1}}

% =====================================================================
%  Tableau de backlog de sprint
% =====================================================================
\newenvironment{backlog}[1]%
  {\begin{table}[H]\centering\caption{#1}%
   \small\begin{tabularx}{\textwidth}{@{}c>{\raggedright\arraybackslash}X c c@{}}
   \toprule
   \textbf{N\textsuperscript{o}} & \textbf{Recit utilisateur} &
   \textbf{Priorite} & \textbf{Estimation} \\ \midrule}%
  {\bottomrule\end{tabularx}\end{table}}
